% SECCIÓN 1. INTRODUCCIÓN
\subsection{¿Qué es el código Morse y por qué es interesante?}
\begin{itemize}
    \item La duración de una raya es tres veces la de un punto
    \item El tiempo que transcurre entre cada raya o punto es igual a la duración de un punto
    \item El espacio entre dos letras tiene la misma longitud que la raya
    \item El espacio entre dos palabras tiene la misma duración que siete puntos
\end{itemize}
\begin{table}[h]
    \centering
    \renewcommand{\arraystretch}{1.2} % Da un poco más de espacio entre las filas
    \begin{tabular}{|c|c|c|c|c|c|}
    \hline
    \textbf{Letra} & \textbf{Código} & \textbf{Letra} & \textbf{Código} & \textbf{Número} & \textbf{Código} \\ \hline
    A & \textbf{$.-$}    & N & \textbf{$-.$}    & 1 & \textbf{$.----$} \\ \hline
    B & \textbf{$-...$}  & O & \textbf{$---$}   & 2 & \textbf{$..---$} \\ \hline
    C & \textbf{$-.-.$}  & P & \textbf{$.--.$}  & 3 & \textbf{$...--$} \\ \hline
    D & \textbf{$-..$}   & Q & \textbf{$--.-$}  & 4 & \textbf{$....-$} \\ \hline
    E & \textbf{$.$}     & R & \textbf{$.-.$}   & 5 & \textbf{$.....$} \\ \hline
    F & \textbf{$..-.$}  & S & \textbf{$...$}   & 6 & \textbf{$-....$} \\ \hline
    G & \textbf{$--.$}   & T & \textbf{$-$}     & 7 & \textbf{$--...$} \\ \hline
    H & \textbf{$....$}  & U & \textbf{$..-$}   & 8 & \textbf{$---..$} \\ \hline
    I & \textbf{$..$}    & V & \textbf{$...-$}  & 9 & \textbf{$----.$} \\ \hline
    J & \textbf{$.---$}  & W & \textbf{$.--$}   & 0 & \textbf{$-----$} \\ \hline
    K & \textbf{$-.-$}   & X & \textbf{$-..-$}  & \multicolumn{2}{c|}{} \\ \cline{1-4}
    L & \textbf{$.-..$}  & Y & \textbf{$-.--$}  & \multicolumn{2}{c|}{} \\ \cline{1-4}
    M & \textbf{$--$}    & Z & \textbf{$--..$}  & \multicolumn{2}{c|}{} \\ \hline
    \end{tabular}
    \caption{Alfabeto internacional y números en Código Morse}
    \label{tab:morse_alfabeto}
\end{table}

\subsection{Objetivo del proyecto}
\begin{itemize}
    \item \textbf{Codificador} - Modo transmisor. Programar el Arduino para que, dada una cadena de texto (por ejemplo, una palabra escrita en el código o enviada desde el PC por el puerto serie), el Arduino la traduzca y la emita haciendo parpadear un LED y sonando un zumbador (buzzer) con los tiempos exactos del código Morse (punto corto, raya larga).
    \item \textbf{Decodificador} - Modo Receptor. Conectar un pulsador (botón) al Arduino. El usuario hará pulsaciones cortas y largas. El programa deberá medir el tiempo que el botón está pulsado para distinguir entre "puntos" y "rayas", agruparlos, y mostrar por el monitor serie (o en una pequeña pantalla LCD) la letra o palabra que se está tecleando.
\end{itemize}